\newpage
\phantomsection
\label{prf:squares-nonnegative}
\LRAProofFor{thm:squares-nonnegative}

\begin{remark*}[Return]
\hyperref[thm:squares-nonnegative]{Return to Theorem}
\end{remark*}

\begin{theorem*}[Squares Are Nonnegative]
For every $a\in F$,
\[
a^2\ge 0.
\]
\end{theorem*}

\begin{proof}[Professional Standard Proof]
\LRAProofBodyStart
Let $a\in F$.  By trichotomy in the ordered field $F$, exactly one of the
following holds:
\[
a>0,\qquad a=0,\qquad -a>0.
\]

If $a=0$, then
\[
a^2=0,
\]
so certainly $a^2\ge 0$.

If $a>0$, then the product of two positive elements is positive.  Hence
\[
a\cdot a>0,
\]
that is,
\[
a^2>0.
\]
Therefore $a^2\ge 0$.

Finally, suppose $-a>0$.  Then again the product of two positive elements is
positive, so
\[
(-a)\cdot(-a)>0.
\]
By the algebraic identity for the product of negatives,
\[
(-a)\cdot(-a)=a\cdot a=a^2.
\]
Thus
\[
a^2>0,
\]
and therefore $a^2\ge 0$.

In every case, $a^2\ge 0$.
\end{proof}

\begin{proof}[Detailed Learning Proof]
\LRAProofBodyStart
We expand the professional proof by following the same sequence of justified moves.

Let $a\in F$.  By trichotomy in the ordered field $F$, exactly one of the
following holds:
\[
a>0,\qquad a=0,\qquad -a>0.
\]

If $a=0$, then
\[
a^2=0,
\]
so certainly $a^2\ge 0$.

If $a>0$, then the product of two positive elements is positive.  Hence
\[
a\cdot a>0,
\]
that is,
\[
a^2>0.
\]
Therefore $a^2\ge 0$.

Finally, suppose $-a>0$.  Then again the product of two positive elements is
positive, so
\[
(-a)\cdot(-a)>0.
\]
By the algebraic identity for the product of negatives,
\[
(-a)\cdot(-a)=a\cdot a=a^2.
\]
Thus
\[
a^2>0,
\]
and therefore $a^2\ge 0$.

In every case, $a^2\ge 0$.
\end{proof}

\begin{remark*}[Proof structure]
The proof splits according to the trichotomy built into the order.  If
$a=0$, the conclusion is immediate.  If $a>0$, positivity of products gives
$a^2>0$.  If $-a>0$, the same positivity argument applies to $(-a)^2$, and
the field identity $(-a)(-a)=a^2$ transfers the conclusion back to $a^2$.
So the square can never land in the negative part of the order.
\end{remark*}

\begin{dependencies}
\begin{itemize}
  \item \hyperref[def:ordered-field]{Definition~\ref*{def:ordered-field}}
        \textnormal{(trichotomy and positivity of products)}
  \item \hyperref[cor:product-of-negatives]{Corollary~\ref*{cor:product-of-negatives}}
\end{itemize}
\end{dependencies}

\clearpage
